\section{Experimental setup}\label{sec:setup}

\subsection{Dataset}\label{sec:data}

We use 3W v2.0.0 \citep{vargas2026} pinned to repository commit \texttt{826fec4} \citep{petrobras3w}: 2,228
instances sampled at 1~Hz, of which 1,119 are real (40 wells), 1,089 simulated and 20 hand-drawn. Class~7
(\emph{scaling in the production choke}) has 36 real instances from six wells (well~24: 19, 22: 8,
23: 4, 21: 2, 6: 2, 1: 1) and 10 hand-drawn ones; there are no simulated class-7 instances. In 31 of the
36 real instances the label never leaves the transient state, so onset is defined by the first transient
or steady label. The pressure downstream of the choke and the downstream temperature are missing in all
real class-7 instances and the choke opening in 78\,\%, which rules out a pressure-drop restriction
index. Under the $+5\,\%$ criterion, 24 of the 36 real instances reach the event (well~24: 17, 22: 3,
6: 2, 1: 1, 21: 1, 23: 0); the median time from a window to the event is 12~h (maximum 147~h).

\subsection{Preprocessing and features}

Signals are averaged to one-minute resolution. Values constant for more than 30~min while the well is
open (frozen sensors) and missing values become mask channels; gaps of up to 10~min are forward-filled;
values with $|x|>10^9$ (sentinel readings) are treated as missing. Variables missing in more than half of
the real instances are dropped, leaving nine channels (upstream choke pressure and temperature, tubing
pressure and temperature, annulus pressure, downhole pressure and temperature, gas-lift line pressure and
the safety-valve state) plus $\widetilde{\RI}$. Robust scaling (median, inter-quartile range) is fitted on
the training files of each fold only. Two data issues of 3W were corrected: (i) 225 file names, covering
885 synthetic instances, repeat across class folders, so instances are keyed by class and file name; and
(ii) hand-drawn pressures are about $10^5$ times smaller than real ones (bar instead of Pa) and are
converted. Windows are formed every 10~min after the first hour, giving 109,547 windows (44,386 real).

\subsection{Splits}

We use six leave-one-well-out folds, one per well with real scaling. In each fold the \emph{test} set
contains all real instances of the held-out scaling well plus a rotating group of wells without scaling,
so every real instance is tested exactly once (24 events, about 248 negative well-days in total). The
\emph{calibration} set contains the next scaling well that has an event (never well~24, which holds 17 of
the 24 events) plus another group of non-scaling wells; it is used for early stopping, hyper-parameter
tuning, alarm thresholds and conformal calibration. The \emph{training} set contains the remaining real
wells and all simulated and hand-drawn instances; synthetic data never enters calibration or test. For the
generalisation test we train on the wells present in 3W v1.1.1 (18 real wells; 4 training events)
and test on wells added in v2.0.0 (wells 22--24, 20 events).

\subsection{Baselines}

All baselines receive the same windows. B0 is a physics rule that thresholds the 60-minute mean of
$\widetilde{\RI}$. B1 (logistic regression) and B2 (gradient-boosted trees \citep{friedman2001}, scikit-learn
\citep{pedregosa2011}) use 63 tabular features at the window end. The deep baselines use per-horizon
binary cross-entropy heads and the same training budget and early stopping as PI-PatchSurv: B3 LSTM
\citep{hochreiter1997}, B4 GRU \citep{cho2014}, B5 CNN-LSTM, B6 TCN \citep{bai2018}, B7 the hybrid
Transformer--LSTM of \citet{lai2026}, B8 PatchTST \citep{nie2023} (our patch encoder without the physics
branch), B8i iTransformer \citep{liu2024itrans}, and B9 an LSTM-autoencoder anomaly detector in the style
of \citet{aranha2024} and \citet{malhotra2016}, trained on normal windows and scored by reconstruction error. B7, B8 and
B9 are our re-implementations of the base-paper architectures under the present protocol.

\subsection{Evaluation metrics}

Positives are rare (about 14\,\% of labelled 24-h windows), so accuracy is not reported. Window-level
metrics per horizon are the area under the precision--recall curve (AUPRC) \citep{saito2015} and under the
ROC curve (AUROC); survival quality is measured by Harrell's C-index \citep{harrell1982}, the Brier score at
24~h and the integrated Brier score over all horizons (IBS) \citep{graf1999}, and calibration by the expected
calibration error (ECE) \citep{guo2017}. Event-level metrics follow deployment: with a threshold fixed on the
calibration wells, an event is \emph{caught} if an alarm fires before $\RI$ reaches $+5\,\%$; we report events
caught out of 24, false alarms per negative well-day, event precision (true over all alarm episodes) and
the lead time. We also report threshold-free recall at budgets of 1 false alarm per well-month and per
well-day. Uncertainty intervals are scored by coverage (PICP) at the 90\,\% target, and $\Delta T$ is the
median absolute error of the predicted time-to-event. Window metrics are averaged over the five folds with
events; pooled AUPRC carries an instance-level bootstrap 95\,\% confidence interval. Significance is assessed
with a Friedman test over folds with the Nemenyi critical difference \citep{demsar2006}, and with a paired
instance-level bootstrap and Wilcoxon signed-rank tests with Holm correction \citep{holm1979}.

\subsection{Implementation details}

All models are implemented in PyTorch \citep{paszke2019} 2.15 (Python 3.14, scikit-learn 1.9, Optuna 5.0,
PHREEQC via phreeqpython). Deep models use $d=64$, AdamW with learning rate $10^{-3}$ and weight decay
$10^{-4}$, batch size 256, 12 epochs of 150 steps and early stopping on the calibration loss; Transformer
layers use four heads, feed-forward width 128 and dropout 0.1. PI-PatchSurv has 128,775 parameters (all
models have fewer than 0.2~M). Every model is run with fixed hyper-parameters and with hyper-parameters
tuned by Optuna \citep{akiba2019} (TPE sampler, seed 0, median pruning, 20 trials per fold) under an equal
budget; tuning is nested in each fold and scored on the calibration-well loss, never on the test well.
Seeds are fixed: PI-PatchSurv and its tuned and \emph{-SI} variants use seeds 0--4, the deep baselines and
the main ablations seeds 0--2, the remaining ablations seed 0. Training ran on a MacBook Air M4 (PyTorch MPS,
float32); seeds 1--2 of the baselines and part of the tuning ran on Kaggle (two NVIDIA T4 GPUs). The
no-patching ablation uses gradient accumulation (micro-batches of 32, effective batch 256) to fit in memory.
An AI assistant (Anthropic Claude) was used to help write and debug the code and to edit the text; all
experiments, results and conclusions were checked by the authors. Code and configuration are publicly
available (see Declarations).
